% 县域农业碳排放核算与减排路径模拟研究
% 安徽省寿县为例
% 统计建模论文 - LaTeX公式版本

\documentclass{article}
\usepackage{ctex}
\usepackage{amsmath}
\usepackage{amssymb}
\usepackage{geometry}
\geometry{left=2.5cm,right=2.5cm,top=2.5cm,bottom=2.5cm}

\title{县域农业碳排放核算与减排路径模拟研究 \\ \small ——以安徽省寿县为例}
\author{}
\date{}

\begin{document}

\maketitle

\section{农业碳排放核算模型}

本文沿用初稿的分项核算思路，将农业碳排放划分为农业投入品隐含碳排放、稻田甲烷排放、农业能源消耗排放和土壤碳汇四部分。

\subsection{总碳排放公式}

\begin{equation}
E_t = E_{input} + E_{CH_4} + E_{energy} - S_{soil}
\end{equation}

其中：
\begin{itemize}
\item $E_t$: 第$t$年农业碳排放净量（$CO_2$当量）
\item $E_{input}$: 农业投入品隐含碳排放
\item $E_{CH_4}$: 稻田甲烷折算的二氧化碳当量
\item $E_{energy}$: 农业能源消耗排放
\item $S_{soil}$: 土壤碳汇量或土壤碳库增量
\end{itemize}

\subsection{农业投入品隐含碳排放}

\begin{equation}
E_{input} = \sum_{m} (X_{mt} \times \delta_m)
\end{equation}

其中：
\begin{itemize}
\item $X_{mt}$: 第$t$年第$m$类农业投入品使用量
\item $\delta_m$: 相应碳排放系数
\end{itemize}

\subsection{稻田甲烷排放}

\begin{equation}
E_{CH_4} = A_t \times \left( E_{FRGS} \times T_s + E_{FnonRGS} \times T_{non} \right) \times GWP_{CH_4}
\end{equation}

其中：
\begin{itemize}
\item $A_t$: 水稻种植面积
\item $E_{FRGS}$: 水稻生长季单位面积甲烷排放量
\item $E_{FnonRGS}$: 非水稻生长季单位面积甲烷排放量
\item $T_s$: 水稻生长季时长
\item $T_{non}$: 非水稻生长季时长
\item $GWP_{CH_4}$: 甲烷全球增温潜势
\end{itemize}

\subsection{农业能源消耗排放}

\begin{equation}
E_{energy} = \sum_{k} (F_{kt} \times EF_k)
\end{equation}

其中：
\begin{itemize}
\item $F_{kt}$: 第$k$类农业能源消费量
\item $EF_k$: 相应排放系数
\end{itemize}

\section{本地化参数修正}

经核验，《气象科学》2025年论文基于寿县2019年11月至2020年10月涡度协方差观测，给出全年$CH_4$通量约$20.52 g \cdot m^{-2}$、水稻生育期约$15.71 g \cdot m^{-2}$；Journal of Integrative Agriculture 2024进一步报告了2019—2021年观测下水稻生长季总$CH_4$通量约$23.9 g \cdot m^{-2}$、非水稻生长季约$3.2 g \cdot m^{-2}$。

\section{K-means聚类识别热点区域}

在乡镇尺度空间识别中，本文以农业碳排放量和碳排放强度为双变量，采用K-means算法对寿县乡镇进行分型。其优化目标函数为：

\begin{equation}
\min \sum_{k=1}^{K} \sum_{x \in C_k} \| x - \mu_k \|^2
\end{equation}

其中：
\begin{itemize}
\item $C_k$: 第$k$类样本集合
\item $\mu_k$: 该类中心
\item $K$: 预设类别数
\end{itemize}

\section{随机森林与SHAP解释模型}

针对农业碳排放影响因素识别，初稿采用随机森林并结合SHAP值解释。

\subsection{随机森林预测形式}

\begin{equation}
f(x) = \frac{1}{B} \sum_{b=1}^{B} T_b(x)
\end{equation}

其中：
\begin{itemize}
\item $T_b(\cdot)$: 第$b$棵回归树的预测结果
\item $B$: 树的数量
\end{itemize}

\subsection{SHAP值一般表达式}

\begin{equation}
\phi_i = \sum_{S \subseteq N \setminus \{i\}} \frac{|S|!(n-|S|-1)!}{n!} \left[ v(S \cup \{i\}) - v(S) \right]
\end{equation}

其中：
\begin{itemize}
\item $\phi_i$: 特征$i$的SHAP值
\item $S$: 特征子集
\item $n$: 特征总数
\item $v(\cdot)$: 模型的输出函数
\end{itemize}

当SHAP绝对值越大，表明该变量对农业碳排放预测的影响越强。

\section{LEAP情景模拟框架}

LEAP模型的核心思想是基于基准年排放清单与活动水平假设，对不同情景下未来排放路径进行递推。

\begin{equation}
E_{ts} = \sum_{j} (AD_{jts} \times EF_{jts})
\end{equation}

其中：
\begin{itemize}
\item $E_{ts}$: 第$t$年情景$s$下的农业碳排放
\item $AD_{jts}$: 第$j$类活动在情景$s$下的活动水平
\item $EF_{jts}$: 对应排放因子或技术校正系数
\end{itemize}

\section{空间格局分析}

如果按空间统计规范继续展开，建议保留全局Moran’s I检验与局部LISA聚类图。其中全局Moran’s I可写为：

\begin{equation}
I = \frac{n \sum_{i,j} w_{ij}(x_i - \bar{x})(x_j - \bar{x})}{\sum_{i,j} w_{ij} \sum_i (x_i - \bar{x})^2}
\end{equation}

其中：
\begin{itemize}
\item $n$: 空间单元数量
\item $w_{ij}$: 空间权重矩阵元素
\item $x_i$: 第$i$个单元的观测值
\item $\bar{x}$: 观测值的均值
\end{itemize}

初稿给出的全局Moran’s I为0.483（$p<0.01$），表明农业碳排放存在显著空间正相关性。

\end{document}
